\chapter{Що насправді каже дофамін}
% full version: chapters/07_dofaalt.tex

У попередньому розділі дофамін був простим: одне число, «краще чи гірше, ніж очікувалося». Це перше наближення, і воно добре працює. Але досліди останніх років показали, що цим «проводом» іде більше, ніж одне число. Цей розділ --- про те, де картина складніша і де науковці поки сперечаються.

\section*{Оптимісти й песимісти}

Дофамінові нейрони не однакові. Одні сильніше реагують на приємні сюрпризи, ніж на неприємні, --- це оптимісти. Інші навпаки --- песимісти. Якщо кожен нейрон учиться за своїм правилом, оптиміст вивчить, якою буде винагорода в найкращому разі, а песиміст --- у найгіршому. Разом вони записують не одне середнє очікування, а весь розкид можливих наслідків, як букмекер, який знає не лише фаворита, а й шанси кожного. Узгодження з дослідом ця ідея має; навіщо мозку весь розкид --- поки гіпотеза.

\pencilfig{7}{\pencilart{7}}{Кожен дофаміновий нейрон тримає свою точку на кривій можливих винагород: разом вони пам'ятають усю криву.}

\section*{Не лише «краще», а й «важливо»}

Частина дофамінових нейронів спалахує і на неприємні, але важливі події: гучний звук, небезпеку. Схоже, що дофаміновою мережею йде не лише знак сюрпризу, а й його вагомість --- сигнал «зверни увагу».

\section*{Два сигнали на одному проводі}

Швидкі сплески навчають. А повільний середній рівень дофаміну, що змінюється за хвилини, схоже, говорить про інше: наскільки багата зараз обстановка. Якщо винагород навколо багато, час дорогий --- і тварина діє швидше. Досліди підтверджують, що рівень дофаміну пов'язаний з енергійністю дій. Інженер назвав би це розділенням сигналів за частотою: високі частоти несуть одне, низькі --- інше.

\section*{Наростання біля мети}

Коли тварина наближається до винагороди, дофамін часто не спалахує, а плавно зростає. Одне пояснення: це саме очікування, сигнал «мета близько, піднаддай». Інше зберігає ідею сюрпризу: здалеку тварина оцінює обстановку розмито, у міру наближення --- дедалі точніше, і кожне уточнення --- маленький приємний сюрприз. Витончений дослід у віртуальному коридорі, де тварину миттєво переносили ближче до мети, свідчить на користь другого: дофамін відповідав сплеском, як і належить сюрпризові.

\section*{Погляд назад}

Є й смілива альтернатива: дофамін відповідає не на питання «що буде?», а на питання «що було причиною винагороди?» --- дивиться в минуле, а не в майбутнє. Досліди, де дві теорії передбачають різне, поки суперечать одне одному. Це жива наукова суперечка.

\section*{Вектор замість числа}

Нарешті, дофамін реагує і на зміну \emph{смаку} винагороди, коли її цінність не змінилася, і допомагає вчитися без жодної винагороди. Схоже, це не одна помилка, а цілий пучок помилок --- по одній на кожну рису того, що сталося.

Висновок розділу: більшість нових теорій не скасовує простої картини, а розширює її. Дофамін --- не один провід, а невеликий джгут із кількох. Лише «погляд назад» сперечається з нею всерйоз.

\fullversion{7}
